\scanpage{14}
\mspage{\handwritten{\unsure{8 b|84|8}}}

dans $G$. Si $C$ est son centralisateur, alors $C$ est représentable et est lisse sur $S$ par \inserted{\handwritten{Exp. $\overline{\underline{\mathrm{XI}}}$}} 5.3. , donc la fonction $s\handwritten{\leadsto}\handwritten{\rho}_n(s)=\dim C_s$ est bien localement constante.

\deleted{Corollaire} La démonstration de \deleted{6}\handwritten{1}.7. est achevée. Nous référerons aux condition envisagées dans \deleted{6}\handwritten{1}.7. b) en disant que dans ce cas, $G$ est de \uline{rang réductif localement constant}.\deleted{\unsure{sur $S$}} Notons:

\uline{Corollaire} \deleted{6}\handwritten{1}.9. Soit $G$ comme dans \deleted{6}\handwritten{1}.7. , et soit $s\handwritten{\in}S$ tel que \inserted{$\handwritten{\rho}_u(s)=0$ i.e.} $\handwritten{\rho}_r(s)=\handwritten{\rho}_{\unsure{n|m}}(s)$,(i.e. les sous-groupes de Cartan de $G_k$ sont des tores, où $k$ est une clôture algébrique de $k(s)$). Alors il existe un voisinage \deleted{\unsure{ouver|ouvert}} ouvert $U$ de $s$ sur lequel $\handwritten{\rho}_r$ et $\handwritten{\rho}_n$ sont constants, en particulier \handwritten{p}our tout $t\handwritten{\in}U$, le rang unipotent $\handwritten{\rho}_u(t)$ de $G_t$ est nul.

En effet, \deleted{\unsure{l'hypothèse s'écrit aussi} $\unsure{\rho}_u(s)=0$, et \unsure{comme}} cela résulte immédiatement de \deleted{6}\handwritten{1}.7. a) et de l'inégalité $\handwritten{\rho}_r(t)\handwritten{\leq}\handwritten{\rho}_{\unsure{n|m}}(t)$ pour tout $t\handwritten{\in}S$.

Notons aussi que nous avons prouvé, en même temps que b), le

\uline{Corollaire} \deleted{6}\handwritten{1}.10. Soit $G$ comme dans \deleted{6}\handwritten{1}.7. , et supposons $G$ de rang réductif localement constant. Considérons le foncteur

\[
\handwritten{\mathscr T}:(\mathrm{Sch})^{\unsure{\mathrm{o}|0|\circ}}\handwritten{\longrightarrow}(\mathrm{Ens})\ ,
\]

tel que pour tout $S'$ sur $S$, on ait

\[
\handwritten{\mathscr T}(S')=\text{ensemble des tores maximaux de }G_{S'}\ .
\]

Alors $\handwritten{\mathscr T}$ est représentable par un préschéma lisse, séparé et de type fini sur $S$.

Il reste à vérifier que $\handwritten{\mathscr T}$ est \deleted{\unsure{séparé}} de type fini sur $S$\handwritten{\unsure{.|?}} La question étant locale pour la topologie fidèlement plate quasi-compacte, on peut supposer que $G$ admet un tore maximal $T$. Mais alors $\handwritten{\mathscr T}$ \deleted{\unsure{se|est}} est \deleted{\unsure{égal à|égal}} \inserted{isomorphe à} $G/\underline{\mathrm{Norm}}_G(T)$ (comparer \inserted{\handwritten{Exp $\overline{\underline{\mathrm{XI}}}$}} 5.3. bis et 5.5. )\handwritten{.} \deleted{\unsure{enfin,|en effet,}} \handwritten{L}e morphisme $G\handwritten{\longrightarrow}\handwritten{\mathscr T}$ défini par $g\handwritten{\leadsto}\mathrm{int}(g)(T)$ étant surjectif , et $G$ quasi-compact sur $S$, il en est de même de $\handwritten{\mathscr T}$ , ce qui achève la démonstration.

% archive sheet number: 13
